\documentclass[11pt]{article}

\usepackage[margin=0.9in]{geometry}
\usepackage{amsmath,amssymb}
\usepackage{booktabs}
\usepackage{enumitem}
\usepackage{listings}
\usepackage{xcolor}
\usepackage{hyperref}
\usepackage{tikz}

\definecolor{backcolour}{rgb}{0.95,0.95,0.92}
\definecolor{codeblue}{rgb}{0.05,0.15,0.55}
\definecolor{codegreen}{rgb}{0.0,0.45,0.0}

\lstset{
    backgroundcolor=\color{backcolour},
    basicstyle=\ttfamily\small,
    keywordstyle=\color{codeblue}\bfseries,
    commentstyle=\color{codegreen},
    stringstyle=\color{purple},
    breaklines=true,
    frame=single,
    showstringspaces=false,
    tabsize=4
}

\title{CPSC 250 \\ Airline Seat Reservation System}
\author{Edward Brash}
\date{}

\begin{document}

\maketitle

\section{Program Goal}

In this example, we design a simple airline seat reservation system.

The goal is to use classes and objects to model seats on an airplane.

Suppose the airplane has five seats. We will create a list containing five \texttt{Seat} objects.

\section{Step 1: Identify Internal Variables}

For each seat reservation, we want to store:

\begin{enumerate}
    \item first name,
    \item last name,
    \item amount paid.
\end{enumerate}

These become the internal variables of the class.

\section{Step 2: Identify Class Methods}

The \texttt{Seat} class should provide methods that allow us to:

\begin{enumerate}
    \item reserve a seat,
    \item unreserve a seat,
    \item check whether the seat is empty,
    \item print seat information.
\end{enumerate}

The user should not directly manipulate the internal variables.

Instead, the user should interact with the object through methods.

\section{The Seat Class}

A reasonable \texttt{Seat} class is:

\begin{lstlisting}[language=Python]
class Seat:

    def __init__(self):
        self.first_name = ""
        self.last_name = ""
        self.amount_paid = 0.0

    def reserve(self, first_name, last_name, amount_paid):
        self.first_name = first_name
        self.last_name = last_name
        self.amount_paid = amount_paid

    def unreserve(self):
        self.first_name = ""
        self.last_name = ""
        self.amount_paid = 0.0

    def is_empty(self):
        return self.first_name == "" and self.last_name == ""

    def print_seat(self):
        print(self.first_name, self.last_name, self.amount_paid)
\end{lstlisting}

\section{Creating Seat Objects}

Suppose:

\begin{lstlisting}[language=Python]
num_seats = 5
available_seats = []
\end{lstlisting}

We can create five seat objects using a loop:

\begin{lstlisting}[language=Python]
for i in range(num_seats):
    available_seats.append(Seat())
\end{lstlisting}

Now \texttt{available\_seats} is a list of objects:

\begin{lstlisting}[language=Python]
[Seat(), Seat(), Seat(), Seat(), Seat()]
\end{lstlisting}

Initially, all seats are empty.

\section{Lists of Objects}

This is an important idea:

\begin{quote}
A list can contain objects, not just primitive values.
\end{quote}

Each element of the list is a separate \texttt{Seat} object.

For example:

\begin{lstlisting}[language=Python]
available_seats[0]
\end{lstlisting}

is the first seat object.

\begin{lstlisting}[language=Python]
available_seats[0].is_empty()
\end{lstlisting}

calls the \texttt{is\_empty()} method for the first seat.

\section{Printing Seat Status}

We can write a function to print the current status of all seats.

\begin{lstlisting}[language=Python]
def print_seats(seats):

    seat_status = []

    for i in range(len(seats)):
        seat_status.append(seats[i].is_empty())

    print("Current seat status")
    print(seat_status)
\end{lstlisting}

Initially this might print:

\begin{lstlisting}
[True, True, True, True, True]
\end{lstlisting}

because all seats are empty.

\section{Printing the Passenger List}

We can also print only the seats that are occupied.

\begin{lstlisting}[language=Python]
def print_passenger_list(seats):

    for i in range(len(seats)):

        if not seats[i].is_empty():

            print(f"{i}: ", end="")
            seats[i].print_seat()
\end{lstlisting}

This demonstrates method calls on objects stored in a list.

\section{Menu Driver}

A menu driver repeatedly asks the user what task to perform.

The options are:

\begin{center}
\begin{tabular}{cl}
\toprule
Command & Meaning \\
\midrule
\texttt{p} & print the seat status \\
\texttt{l} & print the passenger list \\
\texttt{r} & reserve a new seat \\
\texttt{q} & quit \\
\bottomrule
\end{tabular}
\end{center}

\section{Menu Loop}

\begin{lstlisting}[language=Python]
command = input("Enter command: ")

while command != "q":

    if command == "p":
        print_seats(available_seats)

    elif command == "l":
        print_passenger_list(available_seats)

    elif command == "r":
        reserve_seat(available_seats)

    else:
        print("Invalid command")

    command = input("Enter command: ")
\end{lstlisting}

\section{Reserving a Seat}

The \texttt{reserve\_seat()} function asks for a seat number and passenger information.

\begin{lstlisting}[language=Python]
def reserve_seat(seats):

    seat_num = int(input("Enter seat number: "))

    if not seats[seat_num].is_empty():

        print("Seat not empty")

    else:

        first_name = input("Enter first name: ")
        last_name = input("Enter last name: ")
        paid = float(input("Enter amount paid: "))

        seats[seat_num].reserve(first_name, last_name, paid)
\end{lstlisting}

\section{Complete Main Program}

\begin{lstlisting}[language=Python]
num_seats = 5
available_seats = []

for i in range(num_seats):
    available_seats.append(Seat())

command = input("Enter command: ")

while command != "q":

    if command == "p":
        print_seats(available_seats)

    elif command == "l":
        print_passenger_list(available_seats)

    elif command == "r":
        reserve_seat(available_seats)

    else:
        print("Invalid command")

    command = input("Enter command: ")
\end{lstlisting}

\section{Key Takeaways}

\begin{itemize}
    \item Classes are useful for modeling real-world objects.
    \item A seat object stores passenger information.
    \item Lists can store objects.
    \item Functions can receive lists of objects as arguments.
    \item Menu-driven programs use loops and branching.
    \item Encapsulation keeps object data organized and controlled.
\end{itemize}

\end{document}
